\section{Build system}\label{sec:buildsystem}

\subsection{Registry}

We use a \texttt{registry} to manage the available switchable components.
The \texttt{registry} is a singleton class, with one instance per base class.
It stores the available derived classes and their corresponding factory functions, allowing for easy instantiation of the desired components at runtime from the name of the derived class.

Along with each base class, we define domain \texttt{registry} and domain \texttt{registrar} classes corresponding to the base class and its constructor parameters.
The domain \texttt{registry} should be used to access the corresponding derived classes, while the domain \texttt{registrar} is responsible for managing the registration process.

Each derived class is responsible for registering itself with the \texttt{registry}.
For the registration to work, each derived class should define a static instance of the corresponding domain \texttt{registrar} class.
Because we wanted the derived classes to live only in source files in a separate library, we had to ensure that the static instances of the domain \texttt{registrar} classes are not optimized away by the linker.
We achieve this by creating an object library from the source files containing the derived classes and linking it to the executables.

\subsection{CMake setup}

We follow the standard setup of a CMake project.
All headers are located in the \texttt{include/contrast} directory, while the source files are located in the \texttt{src} directory.

The source files are organized into subdirectories based on the library they belong to.
\begin{itemize}
\item \texttt{contrast\_core} is the main library that contains the core functionality of the tool and implementations of the base classes.
\item \texttt{contrast\_common\_derived} is an object library that contains implementations of the derived classes, such as concrete comparison algorithms and printers.
\item \texttt{contrast\_clangtype} contains implementations of abstractions for Clang AST nodes.
\end{itemize}

The \texttt{apps} directory contains the source files for the executables.
We provide a set of main files for different use cases, corresponding to smaller parts of the tool, such as finding crystals (details in \autoref{sec:crystalbleu}), computing similarity scores, and generating reports.

The only part of the project that depends on Clang AST nodes is the \texttt{contrast\_clangtype} library.
The type descriptors for Clang AST nodes are defined in the \texttt{clangast} library inside the \texttt{comparatree} project, which is linked to the \texttt{contrast\_clangtype} library.

If we want \texttt{ContrAST} to work with a different set of original type descriptors, we create a new library with the corresponding abstractions (or stick with \texttt{identity} provided by \texttt{contrast\_common\_derived}).
We then create new CMake targets for the apps, linking them to the new library instead of \texttt{contrast\_clangtype}.
It is important to create the new targets instead of linking the library to the existing ones, as we need to keep the type descriptors from different languages separate to avoid conflicts.

The algorithms and printers from \texttt{contrast\_common\_derived} do not depend on the AST labels, so they can be used with any AST representation.
Provided that we include the abstraction type descriptors, we can also create implementations for algorithms and printers that depend on the concrete AST labels.
This might be useful for algorithms that specifically target certain AST nodes, such as function declarations.
